\section{Introduction et objectifs}

\begin{frame}{Pourquoi la plongée ?}
  \begin{itemize}
    \item La mer : un monde \textbf{merveilleux et mystérieux}.
    \item Découvrir de nouvelles sensations et un \textbf{nouveau mode de communication}.
    \item Apprendre à \textbf{maîtriser son corps} dans un milieu liquide.
    \item La plongée reste un \alert{plaisir} à condition de respecter quelques principes.
  \end{itemize}
\end{frame}

\begin{frame}{Rôle de la formation théorique}
  \begin{itemize}
    \item Acquérir des \textbf{connaissances pratiques et théoriques} adaptées au N1.
    \item Comprendre les \textbf{phénomènes physiques} qui agissent sur le plongeur.
    \item Connaître et appliquer les \textbf{règles de sécurité}.
    \item Savoir reconnaître les \textbf{dangers du milieu} et adopter la bonne attitude.
  \end{itemize}
\end{frame}

\begin{frame}{Objectifs du cours}
  À l’issue de ce cours, le plongeur Niveau 1 devra :
  \begin{itemize}
    \item Situer son niveau dans l’\textbf{organisation fédérale} (FFESSM/CMAS).
    \item Expliquer simplement la \textbf{pression} et la \textbf{loi de Boyle-Mariotte}.
    \item Comprendre la \textbf{flottabilité} et l’utilisation du lest et du gilet.
    \item Identifier les principaux \textbf{accidents barotraumatiques} et leurs préventions.
    \item Comprendre la \textbf{surpression pulmonaire} et l’\textbf{accident de décompression}.
    \item Connaître les \textbf{dangers du milieu} et les \textbf{règles de sécurité}.
  \end{itemize}
\end{frame}

\begin{frame}{Présentation du CST}
  \begin{itemize}
    \item Club associatif affilié à la \textbf{FFESSM}.
    \item Encadrement : E1, E2, E3 / MF1, E4 / MF2.
    \item Formation pratique et théorique vers les \textbf{niveaux fédéraux}.
    \item Sorties, entraînements, et \textbf{esprit de club}.
  \end{itemize}
\end{frame}
